\part{Des échantillons décodés à une photographie interprétée}
\section{Décodage, développement et rendu : trois contrats}
\subsection{Décodage du fichier}
Le décodage identifie le format, lit les structures, localise
les charges utiles et reconstruit les échantillons déclarés.
Pour un RAW mosaïqué, la sortie peut rester une seule grille
de codes par photosite. Pour un TIFF RGB, elle peut être un
tableau de composantes déjà rendues. Les deux sorties n'ont
pas la même signification, même si elles utilisent uint16.

\subsection{Développement du RAW}
Le développement utilise les mesures et les calibrations pour
produire une image interprétable : niveaux, corrections,
balance des blancs, reconstruction colorée, transformation
vers un espace de travail et éventuellement rendu de base.
LibRaw sépare notamment ouverture, dépaquetage et étapes de
traitement dans son API ; lire \texttt{open\_file} ou réussir
\texttt{unpack} n'est pas obtenir automatiquement un rendu
final~\cite{libraw}.

\subsection{Rendu et affichage}
Le rendu décide comment représenter une scène dans les
contraintes du support. Il comprend éventuellement courbes,
compression des hautes lumières, choix esthétiques, réduction
du gamut et conversion vers l'écran. Il n'existe pas un
seul rendu mathématiquement obligatoire pour toute scène.
En revanche, une interprétation incorrecte du noir, de la
matrice ou de la fonction de transfert est un défaut
technique, pas simplement un style artistique.

\resultat{Deux logiciels peuvent décoder les mêmes mesures puis
produire des rendus différents. Pour identifier un bug, il faut
comparer les étapes, pas seulement les images finales.}

\section{Un pipeline RAW conceptuel détaillé}
\begin{center}
\begin{tikzpicture}[node distance=.3cm,
box/.style={draw=iprblue,rounded corners,text width=12.5cm,align=center,minimum height=.66cm}]
\node[box] (f) {Fichier : structures, métadonnées, aperçus et données RAW};
\node[box,below=of f] (r) {Décompression et dépaquetage : grille native};
\node[box,below=of r] (n) {Niveaux, linéarisation éventuelle, pixels défectueux et corrections capteur};
\node[box,below=of n] (w) {Gains de balance des blancs et reconstruction des canaux};
\node[box,below=of w] (p) {Transformation caméra, profil, tables et rendu tonal de base};
\node[box,below=of p] (e) {Réglages, retouches, géométrie et traitement du détail};
\node[box,below=of e] (o) {Rendu de sortie, gamut, fonction de transfert, quantification et fichier};
\draw[-{Latex}] (f)--(r);\draw[-{Latex}] (r)--(n);\draw[-{Latex}] (n)--(w);
\draw[-{Latex}] (w)--(p);\draw[-{Latex}] (p)--(e);\draw[-{Latex}] (e)--(o);
\end{tikzpicture}
\end{center}
Cet ordre est pédagogique, pas universel. Certains débruitages
opèrent avant dématriçage ; certaines corrections optiques
après reconstruction ; des profils imposent un domaine
et un ordre précis. Toute permutation doit être justifiée.

\subsection{Dépaquetage et précision native}
Un codec peut reconstruire un flux de codes 12 ou 14 bits,
qu'il faut ensuite ranger dans un tampon. Le dépaquetage
traduit leur disposition binaire, sans décider de leur couleur.
Un ordre des bits incorrect peut être masqué sur des faibles
valeurs et devenir visible sur les plus fortes. Les tests
doivent inclure zéro, maximum et valeurs traversant les
frontières d'octets.

\subsection{Linéarisation, noir et blanc}
Certains fichiers décrivent une table de linéarisation ou
des niveaux qui ne se résument pas à deux constantes. Une
normalisation doit respecter les données déclarées. Dans
le cas simple, on utilise l'équation~\eqref{eq:normalisation}.
Si les noirs varient entre canaux, lignes ou colonnes,
écrire $B_c(x,y)$ plutôt qu'une constante universelle.

Les valeurs négatives après soustraction du noir peuvent
représenter des fluctuations normales autour de zéro.
Les couper immédiatement introduit un biais positif dans
les statistiques des ombres. À l'inverse, certaines étapes
comme des tables définies seulement sur $[0,1]$ exigent une
politique explicite de bornage. Les deux impératifs doivent
être conciliés, pas ignorés.

\subsection{Correction des défauts et optique}
Les photosites chauds ou morts peuvent être remplacés par
une estimation locale, éventuellement adaptée au CFA.
Une correction de vignettage multiplie spatialement les
valeurs ; elle amplifie aussi le bruit des zones sombres.
Une correction géométrique rééchantillonne la grille ; elle
change donc les valeurs et les correspondances pixel à pixel.
Une correction d'aberration chromatique peut déformer
différemment les canaux. Cela rend une MAE sans recalage
géométrique difficile à interpréter.

\subsection{Balance des blancs}
Un modèle diagonal courant est
\begin{equation}
\begin{pmatrix}R'\\G'\\B'\end{pmatrix}=
\begin{pmatrix}g_R&0&0\\0&g_G&0\\0&0&g_B\end{pmatrix}
\begin{pmatrix}R\\G\\B\end{pmatrix}.
\end{equation}
Les gains tentent de rendre neutre une référence éclairée
par une lumière donnée dans un contexte de calibration.
Multiplier les gains par une constante commune change
aussi la luminosité ; leur normalisation importe.
Un Kelvin affiché est une description de chromaticité ou
un paramètre d'interface, pas les trois gains eux-mêmes.
Deux moteurs peuvent afficher des Kelvin différents à
partir de gains proches, ou l'inverse, suivant leurs modèles.

Sous lumière mixte, une seule balance globale ne peut
neutraliser toutes les régions. Sous un projecteur à spectre
discontinu, « corriger la température » ne recrée pas les
longueurs d'onde absentes. La nuance fournit un degré
supplémentaire, mais pas une reconstruction spectrale complète.

\subsection{Dématriçage : une estimation, pas une lecture directe}
Sur une mosaïque RGGB, on peut estimer un vert à la position
d'un rouge par la moyenne des quatre voisins verts :
\begin{equation}
\widehat G(x,y)=\tfrac14[G(x-1,y)+G(x+1,y)+G(x,y-1)+G(x,y+1)].
\end{equation}
C'est une approximation bilinéaire simple. Elle lisse des
contours et peut produire des fausses couleurs. Des méthodes
adaptatives choisissent une direction ou exploitent la
corrélation des différences de couleur. Les méthodes
apprises proposent d'autres compromis, mais exigent des
données d'entraînement et une validation de généralisation.

La comparaison de deux dématriçages doit séparer détail,
bruit, moiré et rendu couleur. Plus net ne signifie pas
nécessairement plus fidèle : une arête peut être un artefact.

\section{La couleur n'est pas contenue dans les nombres seuls}
\subsection{Du spectre aux composantes}
Une réponse de caméra peut s'écrire
\begin{equation}
c_k=\int S(\lambda)\rho(\lambda)s_k(\lambda)\,d\lambda.
\end{equation}
Les fonctions $s_k$ sont les sensibilités spectrales effectives
des canaux. Les tristimulus d'un observateur standard
s'obtiennent avec d'autres fonctions de pondération, par
exemple $\overline x,\overline y,\overline z$ pour XYZ.
Trois nombres de caméra ne sont donc pas directement les
trois nombres d'un espace RGB d'affichage.

Deux spectres distincts peuvent donner les mêmes trois
réponses : c'est une forme de métamérisme. Une matrice
3$\times$3 peut approximer une transformation sur un domaine
donné, mais ne reconstitue pas tout le spectre. Une charte
photographiée sous un illuminant fournit une calibration
utile ; elle ne garantit pas la perfection sous tous les
éclairages et pour toutes les matières.

\subsection{Matrices et changement de base}
Dans un modèle linéaire :
\begin{equation}
\begin{pmatrix}X\\Y\\Z\end{pmatrix}=M
\begin{pmatrix}R_c\\G_c\\B_c\end{pmatrix}.
\end{equation}
Chaque colonne indique la contribution d'un canal caméra.
La direction de la matrice est essentielle : caméra vers
XYZ et XYZ vers caméra sont deux transformations inverses,
pas des noms interchangeables. Certaines calibrations
incluent des conventions de balance ou d'adaptation.
Appliquer deux fois une correction qui était déjà incorporée
produit une dérive.

Un espace RGB est défini par des primaires, un point blanc
et une convention de codage. Le passage entre primaires
linéaires peut être une matrice ; le passage entre codes
non linéaires exige d'abord l'inversion de leurs fonctions
de transfert. Une matrice appliquée directement à des
codes sRGB ne réalise généralement pas le bon changement
de base physique.

\subsection{Point blanc et adaptation chromatique}
D50 et D65 sont des blancs de référence, pas les températures
mesurées d'une photo particulière. L'adaptation chromatique
modélise une partie de l'adaptation visuelle lorsque l'on
change de blanc. Un modèle de type von Kries transforme
vers un espace de réponses, multiplie par des gains, puis
revient : $M^{-1}\operatorname{diag}(d)M$.
Bradford est un choix de matrice de ce type, développé dans
le cours 03. Confondre D50 et D65 peut déplacer les couleurs
même lorsque toutes les multiplications sont numériquement
correctes.

\section{Fonctions de transfert, gamma et HDR}
\subsection{Linéaire veut dire linéaire par rapport à quoi ?}
Dans un domaine linéaire par rapport à la lumière, doubler
un signal non saturé double la grandeur représentée. Une
valeur sRGB de 0,5 n'est pas la moitié de la lumière de 1.
La fonction d'encodage sRGB normalisée est~\cite{srgbimages}
\begin{equation}
E(L)=\begin{cases}
12{,}92L,&L\le0{,}0031308,\\
1{,}055L^{1/2{,}4}-0{,}055,&L>0{,}0031308.
\end{cases}
\end{equation}
Son inverse usuel est
\begin{equation}
L(E)=\begin{cases}
E/12{,}92,&E\le0{,}04045,\\
((E+0{,}055)/1{,}055)^{2{,}4},&E>0{,}04045.
\end{cases}
\end{equation}
Ainsi $E=0{,}5$ donne $L\approx0{,}21404$ et
$L=0{,}18$ donne $E\approx0{,}46136$. « Gamma 2,2 » est une
approximation fréquente, pas la fonction sRGB exacte.

\subsection{Pourquoi l'ordre des opérations compte}
La moyenne linéaire d'un noir et d'un blanc vaut $L=0{,}5$,
soit $E\approx0{,}73536$ après encodage. Faire la moyenne
des codes donne $E=0{,}5$, soit seulement $L\approx0{,}21404$.
On comprend ainsi pourquoi une interpolation dans le mauvais
domaine peut produire des transitions trop sombres.

Une correction d'exposition physique utilise
$L'=2^eL$, et non généralement $E'=2^eE$.
À $+1$ EV, un code sRGB 0,5 devient environ 0,68584
tant que le signal reste dans le domaine autorisé ; il ne
doit pas automatiquement devenir 1.

Quantifier avant l'encodage puis après lui n'est pas non
plus équivalent :
\[
Q(E(L))\ne Q(E(Q(L))).
\]
Une quantification linéaire prématurée peut supprimer des
faibles valeurs que l'encodage aurait réparties sur davantage
de codes. Ce problème a motivé des contrôles spécifiques
de sortie sRGB dans IPR ; ce n'est pas une propriété
inévitable de tout TIFF.

\subsection{OETF, EOTF et courbe artistique}
Une OETF traduit un signal optique vers un signal électrique
dans une chaîne donnée ; une EOTF décrit le comportement
d'un affichage à partir de ses codes. Une transformation
globale de la scène à l'écran peut aussi inclure une OOTF.
La courbe artistique ajuste le rendu tonal. Ce sont des
objets différents, même si plusieurs sont non linéaires.
Appliquer une courbe de profil ne dispense pas de décrire
l'encodage final ; appliquer deux encodages sRGB assombrit
ou éclaircit fortement le résultat selon la direction.

\subsection{HDR n'est pas synonyme de 16 bits}
HDR peut décrire une scène, un espace de calcul étendu ou
une diffusion pour un écran à grande luminance. Un TIFF
16 bits peut rester SDR ; un fichier à 10 bits correctement
décrit peut porter un rendu HDR. La recommandation
BT.2100 encadre notamment PQ et HLG~\cite{hdr}.

PQ relie les codes à une luminance d'affichage absolue selon
une fonction perceptuelle normalisée. HLG emploie une
construction différente, avec un comportement système lié
aux conditions d'affichage. Il faut conserver la fonction
de transfert, les primaires, la plage de codes et les
métadonnées pertinentes. Ouvrir les codes HDR comme sRGB
donne une image incohérente, même sans erreur de décompression.

Une \emph{gain map} peut décrire comment transformer une
base SDR vers un rendu à plus grande dynamique. Elle n'est
ni nécessairement une deuxième photographie indépendante,
ni une preuve que le capteur avait conservé toutes les
hautes lumières. Un moteur doit savoir s'il la prend en
charge ; l'ignorer peut laisser un rendu SDR utilisable,
mais différent du rendu voulu sur écran HDR.

\section{Profils ICC, DCP et LCP : des rôles distincts}
\subsection{ICC : décrire une relation colorimétrique}
Un profil ICC décrit une transformation entre un espace
de périphérique ou de travail et un espace de connexion
normalisé, suivant sa classe et sa structure. Des profils
peuvent utiliser matrices/courbes ou tables. Le PCS des
profils ICC classiques utilise notamment une référence D50 ;
les intentions de rendu guident le traitement des couleurs
non reproductibles~\cite{iccimages}.

\emph{Assigner} un profil change l'interprétation des nombres
sans forcément les changer. \emph{Convertir} cherche à
préserver l'apparence ou une relation colorimétrique en
calculant de nouveaux nombres. Assigner sRGB à des valeurs
Adobe RGB ne remplace donc pas une conversion Adobe RGB
vers sRGB. Un profil écran est destiné à la sortie du
système d'affichage, pas à devenir automatiquement le
profil incorporé à chaque photo exportée.

\subsection{DCP : caractériser et rendre des valeurs caméra}
Un DCP peut comporter matrices, illuminants, tables de
teinte/saturation, look, courbe et offset d'exposition.
Ces éléments n'ont pas tous un rôle purement colorimétrique :
certains participent au rendu de base. Leur domaine et leur
ordre sont essentiels~\cite{dng}. Une table décrite dans
des primaires de travail précises ne doit pas être évaluée
comme si elle était une table sRGB quelconque.

Une interpolation entre deux illuminants peut utiliser
des températures réciproques. Par exemple,
\begin{equation}
w=\operatorname{clip}\left(
\frac{1/T-1/T_1}{1/T_2-1/T_1},0,1\right),\qquad
M=(1-w)M_1+wM_2.
\end{equation}
C'est une règle de mélange d'un modèle, pas une interpolation
exacte de tous les spectres lumineux possibles.

Pour une courbe RGB, appliquer $T$ à chaque canal peut
changer la position relative du canal médian et donc la
teinte. Une alternative applique la courbe aux minimum et
maximum puis interpole le médian. DCamProf explique
l'intérêt de cette stabilisation et les interactions
avec les tables de rendu~\cite{dcamprofimages}. Elle ne
garantit ni saturation constante ni fidélité parfaite
après réduction du gamut. Le cours ne recommande pas
une substitution aveugle au moteur de production.

\subsection{LCP et corrections d'objectif}
Un profil d'objectif peut décrire distorsion, vignettage
et aberrations suivant focale, ouverture ou distance.
Son action peut être géométrique et photométrique. Il ne
remplace pas le profil couleur du capteur. Une correction
de distorsion déplace les coordonnées : pour des tests
pixel à pixel, son activation doit être harmonisée.

\subsection{Légalité et portabilité}
Une spécification de format ouverte et les données de
profils distribuées par un éditeur sont deux objets
différents. Le fait de pouvoir lire un profil ne donne
pas automatiquement le droit de le redistribuer avec
un logiciel. Les licences des bibliothèques, des profils,
des chartes et des jeux d'images doivent être examinées
séparément. Cette distinction n'est pas un avis juridique.

\section{Les métadonnées : plusieurs familles et plusieurs vérités}
\subsection{Inventaire utile}
\begin{longtable}{p{3.1cm}p{6.5cm}p{5cm}}
\toprule Famille & Exemples & Conséquence \\ \midrule
Structure & Dimensions, profondeur, compression, CFA & Nécessaire au décodage. \\
Prise de vue & Temps, ouverture, ISO, focale, date & Contexte, pas nécessairement une transformation. \\
Calibration & Noir, blanc, gains, matrices & Change l'interprétation du signal. \\
Présentation & Orientation, profil, miniature & Change le rendu ou sa géométrie. \\
Éditorial & Auteur, légende, mots-clés, droits & Organisation et diffusion. \\
Géographie & Latitude, longitude, altitude & Information sensible à protéger. \\
Développement & Exposition, courbe, paramètres de profil & Recette, parfois déjà incorporée aux pixels. \\
Gestion & Note, drapeau, libellé, identifiant & Peut vivre dans un catalogue ou sidecar. \\
\bottomrule
\end{longtable}

\subsection{Exif, MakerNotes, IPTC et XMP}
Exif fournit des conventions de métadonnées photographiques
et des structures associées aux images ; la page CIPA
recense les versions de la norme, dont Exif~3.1 publiée
en 2026~\cite{exif}. Un champ ISO peut figurer en plusieurs
endroits ou avec plusieurs interprétations historiques.
Il faut afficher les groupes de tags, pas seulement leurs
noms courts.

Les MakerNotes contiennent des données propres au
constructeur. Elles peuvent avoir leurs propres offsets,
versions et conventions. ExifTool documente les champs
qu'il sait interpréter ; cette documentation est une
source d'implémentation précieuse, pas une norme constructeur
universelle~\cite{exiftoolimages}. Déplacer un bloc sans
ajuster ses offsets peut casser sa lecture.

IPTC couvre notamment les usages éditoriaux ; XMP propose
un modèle extensible sérialisé couramment en RDF/XML,
embarqué ou externe~\cite{xmpimages,iptcimages}. Un catalogue
peut aussi garder des réglages sans les écrire dans le RAW.
« Je vois une étoile dans Lightroom » ne prouve donc pas
que cette étoile est présente dans le fichier transmis.

\subsection{Recette et pixels déjà développés}
Un TIFF exporté depuis Lightroom peut contenir des
métadonnées de développement ainsi que des pixels où le
développement est déjà effectué. Réappliquer les réglages
XMP sur ces pixels applique potentiellement l'effet deux
fois. Un moteur doit identifier la nature de l'entrée,
pas interpréter tout paramètre présent comme une instruction
de rendu à exécuter.

Un sidecar \texttt{.ipr} ou \texttt{.xmp} est une recette
non destructive seulement si l'original est préservé et
si les paramètres, versions et ressources nécessaires
restent disponibles. Pour reproduire un rendu, conserver
aussi le profil exact, la version moteur et les conventions.

\subsection{Données absentes, inconnues ou contradictoires}
Une valeur absente n'est pas automatiquement zéro.
« GPS absent » ne signifie pas latitude 0, longitude 0.
« Profil absent » ne prouve pas sRGB. Une température
calculée par un logiciel n'est pas forcément la valeur
inscrite par le boîtier. La stratégie correcte est de
conserver provenance, statut connu/inconnu et hypothèses
de repli, avec un avertissement si la comparaison dépend
de cette hypothèse.

\section{Pourquoi le même RAW change d'apparence selon le logiciel}
Les différences peuvent provenir du dépaquetage, du noir,
de la normalisation, de la WB, du dématriçage, des matrices,
des tables, de la courbe de base, de la récupération des
hautes lumières, des corrections optiques, du bruit,
du gamut mapping ou de l'affichage. Une miniature Canon
est un rendu déjà produit par Canon, pas une vue directe
des valeurs CFA. La comparaison à l'œil de cette miniature
et d'un export ne permet pas de localiser l'étape différente.

Le bon protocole compare d'abord les mêmes fichiers,
mêmes profils, mêmes réglages, mêmes fonctions de transfert
et même géométrie, puis isole les étages. Un profil nommé
« Neutral » n'est pas forcément une matrice sans courbe.
« Tous les curseurs à zéro » ne signifie pas « aucun
traitement ». Cette distinction est centrale pour la
validation de la base colorimétrique d'IPR.
